\section{Common reductions and the interval target}
\label{sec:preliminaries}
We use $\binom ab=0$ when $b<0$ or $b>a$, for $a\ge0$.
Unless expressly indicated otherwise, differences in numerical
inequalities are ordinary integer or rational differences. Rational
weights in the proof are counting coefficients, and impose no condition
on the characteristic of $k$.

\subsection{Orders, completion, and scalar extension}
For $s\in S$, let $\ord(s)$ be the largest length of an expression of
$s$ as a sum of minimal generators. If $m$ is the multiplicity, the
Ap\'ery set $\Ap(S,m)=\{s\in S:s-m\notin S\}$ contains one element
in each residue class modulo $m$. Every $s\in S$ has a unique expression
$s=b+\ell m$ with $b\in\Ap(S,m)$ and $\ell\ge0$.

\begin{proposition}\label{prop:bridge}
The following operations preserve the objects needed in
Theorem~\ref{thm:main}.
\begin{enumerate}
\item The associated graded rings of $k[t^S]$ and $k[[t^S]]$ at their
positive monomial ideals are naturally isomorphic, preserving the
initial form of every minimal generator.
\item For every field extension $K/k$,
$\beta_i(G_S;k)=\beta_i(G_S\otimes_k K;K)$ for all $i$.
\item If $w\le m-2$ and $e\ge3$, there are integers $d\ge1$,
$u\ge e-1$, and $0<a_1<\cdots<a_q<u$, with $q=e-2$, such that
\begin{gather}\label{eq:normalization}
 S=\langle m,m+da_1,\ldots,m+da_q,m+du\rangle,\\
 \gcd(a_1,\ldots,a_q,u)=\gcd(m,d)=1,\qquad m\ge du+2.\nonumber
\end{gather}
Here $du=w$ is the physical width.
\end{enumerate}
\end{proposition}
\begin{proof}
The $n$th power of the positive monomial ideal consists exactly of the
elements supported on exponents $s$ with $\ord(s)\ge n$. This holds
also for formal series: assign each such exponent to one of the finitely
many length-$n$ products that divides it, and group its coefficient into
the corresponding formal series. The degree-$n$ quotient consequently
has basis the exponents of order exactly $n$. This set is finite, since
$s=b+\ell m$ of order $n$ has $\ell\le n$. Both associated graded
rings have these bases, and multiplication agrees: a product survives
exactly when orders add. This proves the first assertion, including
compatibility with the polynomial-ring action.

Tensor a minimal graded free resolution with $K$. Flatness preserves
exactness, and all differential entries remain in the homogeneous
maximal ideal, so minimality and ranks are preserved. The preceding
monomial bases also identify the scalar extension of the tangent cone.

For the last assertion, take
$d=\gcd(g_2-m,\ldots,g_e-m)$ and $u=(g_e-m)/d$.
Dividing the internal offsets by $d$ gives the $a_j$. Their common gcd
with $u$ is one, while $\gcd(m,d)=1$ follows from $\gcd(S)=1$.
There are $q$ distinct internal positive integers below $u$, giving
$u\ge q+1=e-1$. The narrow-width hypothesis gives $m\ge du+2$.
\end{proof}

In particular, an infinite extension may be used for a general linear
cut and then removed. We never assume that the whole ring $G_S$ is
finite dimensional over $k$; finiteness is used for its graded pieces,
Artinian quotients, and Koszul homology.

\subsection{Exact target and uniform capacity}
Suppose $m\ge2$, and put
\[
 v=\min\{M-m,m-1\},\qquad
 D=v-((m-1)\bmod v),\quad 1\le D\le v.
\]
The minimal generators of $\widetilde S$ are precisely
$m,m+1,\ldots,m+v$. If $M<2m$, all listed generators are
indecomposable. Otherwise $m,\ldots,2m-1$ generate every integer
at least $m$, by induction after subtraction of $m$; generators above
$2m-1$ are redundant.

\begin{proposition}\label{prop:capacity}
For every $i\ge1$, the target Betti number equals
\begin{equation}\label{eq:capacity}
 F_i(m,v)=i\binom v{i+1}+
 \begin{cases}
 (D+1-i)\binom v{i-1},&i\le D,\\
 (i-D)\binom vi,&i>D.
 \end{cases}
\end{equation}
For $1\le i\le v$, it satisfies
\begin{equation}\label{eq:mincapacity}
 F_i(m,v)\ge F_i^{\min}(v):=
 i\binom v{i+1}+\min\left\{\binom v{i-1},\binom vi\right\}.
\end{equation}
For fixed $i$, $F_i^{\min}(v)$ is nondecreasing on $v\ge i$.
If $v=m-1$, then
\begin{equation}\label{eq:ordinary-capacity}
 F_i(m,m-1)=i\binom m{i+1}.
\end{equation}
\end{proposition}
\begin{proof}
The arithmetic-sequence Betti formula \cite[Theorem~4.1]{APFormula}
and equality of semigroup-ring and tangent-cone Betti numbers
\cite[Proposition~2.5]{APCone} apply to the minimal interval generators.
Writing $m=av+b$ with $1\le b\le v$ gives $D=v-b+1$ and yields
\eqref{eq:capacity}, with the quotient-ring homological indexing.
On the branch $D\ge i$ the correction is at least $\binom v{i-1}$;
on $D<i$ it is at least $\binom vi$. This proves
\eqref{eq:mincapacity}, including $i=1$ and $i=v=1$.
Each binomial coefficient increases with $v$, proving monotonicity.
For $v=m-1$ one has $D=v$; Pascal's identity gives
\eqref{eq:ordinary-capacity}.
\end{proof}

Some useful consequences, with their necessary width ranges, are
\begin{align}\label{eq:low-capacities}
 F_1(m,v)&\ge\binom v2+1 &&(v\ge1),\nonumber\\
 F_2(m,v)&\ge2\binom v3+v &&(v\ge3),\\
 F_3(m,v)&\ge3\binom v4+\binom v2 &&(v\ge5),\nonumber\\
 F_4(m,v)&\ge4\binom v5 &&(v\ge4).\nonumber
\end{align}
We do not extend the third-row estimate to $v=4$.
In the narrow case $v=du\ge u$. Consequently, whenever $u\ge i$,
an upper bound $\beta_i(G_S;k)\le F_i^{\min}(u)$ pays for the
actual target by monotonicity. At the smallest normalized widths, it is
sometimes more efficient to use the physical width directly or to
identify source and target.

\subsection{The width bound for the local ring}
\begin{proof}[Proof of Corollary~\ref{cor:width-bound}]
Put $\widehat P_S=k[[X_1,\ldots,X_e]]$. The standard filtered
Koszul comparison gives
$\beta_i(R_S;k)\le\beta_i(G_S;k)$. Indeed, filter the Koszul
complex on the images of $X_1,\ldots,X_e$ by powers of the maximal
ideal, shifting the filtration by the exterior degree. Its associated
graded complex is the Koszul complex of the initial forms on $G_S$.
The resulting convergent spectral sequence starts with
$\Tor^{P_S}(G_S,k)$ and abuts to the associated graded of
$\Tor^{\widehat P_S}(R_S,k)$. Completeness and the finite-dimensional
first-page homology give the dimension inequality in each homological
degree. This argument works over every field.

If $w=0$, then $S=\N$ and all positive Betti numbers vanish.
Otherwise use Theorem~\ref{thm:main} and put
$v=\min\{w,m-1\}$. In \eqref{eq:capacity}, the correction for
$i\le D$ is at most
\[
 (v+1-i)\binom v{i-1}=i\binom vi;
\]
for $i>D$ it is also at most $i\binom vi$. Hence
\[
 \beta_i(G_{\widetilde S};k)=F_i(m,v)
 \le i\binom{v+1}{i+1}\le i\binom{w+1}{i+1}.
\]
For the stated ordinary semigroup, $m=w+1$ and
\eqref{eq:ordinary-capacity} is equality in the width bound.
The arithmetic-sequence semigroup ring and its tangent cone have the
same total Betti numbers \cite{APFormula,APCone}, giving equality also
for $R_S$.
\end{proof}
